\documentclass[11pt]{article}

% --- UNIVERSAL PREAMBLE BLOCK ---
% Set 4:5 Aspect Ratio (Standard LinkedIn Portrait Carousel 1080x1350 equivalent)
\usepackage[paperwidth=8in, paperheight=10in, top=1.2in, bottom=1.2in, left=1in, right=1in]{geometry}
\usepackage{fontspec}
\usepackage[english, bidi=basic, provide=*]{babel}

\babelprovide[import, onchar=ids fonts]{english}

% Set premium typography using Noto Sans system fonts
\babelfont{rm}{Noto Sans}
\babelfont{sf}{Noto Sans}
\babelfont{tt}{Noto Sans Mono}

% --- PACKAGES FOR MODERN DESIGN & MATHS ---
\usepackage{amsmath}
\usepackage{amssymb}
\usepackage{xcolor}
\usepackage{listings}
\usepackage{mdframed}
\usepackage{tikz}
\usetikzlibrary{shapes.geometric, arrows.meta, positioning, shadows.blur}
\usepackage{tcolorbox}
\tcbuselibrary{listings, skins, breakable}
\usepackage{fancyhdr}
\usepackage{enumitem}

% --- DESIGN TOKENS (Stripe & Linear Inspired) ---
\definecolor{brandnavy}{HTML}{0A2540}    % Deep Primary Navy
\definecolor{brandgray}{HTML}{425466}    % Slate Gray
\definecolor{brandlight}{HTML}{F8F9FA}   % Light Soft Gray
\definecolor{brandblue}{HTML}{635BFF}    % Accent Classic Blue
\definecolor{brandgreen}{HTML}{00D4B2}   % Success/Muted Green
\definecolor{brandorange}{HTML}{F25C05}  % Accent Orange
\definecolor{brandred}{HTML}{E5484D}     % Warning Red

% Premium Dark Theme for Code Blocks
\definecolor{codebackground}{HTML}{141419}
\definecolor{codekeyword}{HTML}{FF7B72}
\definecolor{codestring}{HTML}{A5D6FF}
\definecolor{codecomment}{HTML}{8B949E}
\definecolor{codetext}{HTML}{E6EDF3}
\definecolor{codeframe}{HTML}{30363D}

% --- ZERO PARAGRAPH INDENT ---
\setlength{\parindent}{0pt}
\setlength{\parskip}{1.2em}

% Custom TColorBox for styled content blocks
\newtcolorbox{infobox}[1][]{
    colback=brandlight, colframe=brandblue, arc=5pt, boxrule=0.8pt,
    left=15pt, right=15pt, top=12pt, bottom=12pt, #1
}

\newtcolorbox{darkbox}[1][]{
    colback=brandnavy, colframe=brandnavy, arc=5pt, boxrule=0.8pt,
    left=15pt, right=15pt, top=12pt, bottom=12pt, 
    coltext=white, #1
}

% Custom TColorBox for Syntax Highlighted Code
\newtcblisting{codebox}[2][]{
    colback=codebackground, colframe=codeframe, arc=5pt, boxrule=1pt,
    left=10pt, right=10pt, top=10pt, bottom=10pt,
    listing only, listing options={
        basicstyle=\ttfamily\footnotesize\color{codetext},
        keywordstyle=\color{codekeyword}\bfseries,
        stringstyle=\color{codestring},
        commentstyle=\color{codecomment}\itshape,
        showstringspaces=false,
        breaklines=true,
        language=#2
    },
    #1
}

% --- HEADER & FOOTER CONFIGURATION ---
\pagestyle{fancy}
\fancyhf{}
\renewcommand{\headrulewidth}{0pt}
\renewcommand{\footrulewidth}{0.5pt}

\fancyhead[L]{\fontsize{9}{11}\selectfont\color{brandgray}\MakeUppercase{MLOps Deep Dive}}
\fancyhead[R]{\fontsize{9}{11}\selectfont\color{brandblue}\textbf{PRODUCTION PARITY}}

\fancyfoot[L]{\fontsize{8}{10}\selectfont\color{brandgray}
Author: Dibayendu Mukherjee \hspace{0.7em}$\cdot$\hspace{0.7em}
\mbox{AI Research Engineer \hspace{0.7em}\textbullet\hspace{0.7em} Systems Engineer}}

\fancyfoot[R]{\fontsize{8}{10}\selectfont\color{brandgray}Slide \thepage}

\begin{document}

% ==========================================================
% SLIDE 1: COVER
% ==========================================================
\thispagestyle{empty}
\vspace*{0.8in}

\begin{center}

    {\fontsize{14}{16}\selectfont\color{brandblue}\textbf{
    FROM NOTEBOOK TO PRODUCTION
    }}

    \vspace{0.8in}

    {\fontsize{40}{48}\selectfont\color{brandnavy}\bfseries
Why my perfect\\[0.15in]
notebook model\\[0.25in]
failed in production.
}

    \vspace{0.6in}

    {\fontsize{18}{24}\selectfont\color{brandgray}
    The painful lesson of silent failures,\\
    mismatched defaults, and why parity\\
    beats perfect metrics.
    }

    \vfill

    \begin{minipage}{0.8\textwidth}
        \centering
        \rule{0.6\linewidth}{0.5pt}

        \vspace{0.2in}

        {\fontsize{12}{14}\selectfont\color{brandnavy}
        \textbf{Dibayendu Mukherjee}
        }

        \vspace{0.05in}

        {\fontsize{10}{12}\selectfont\color{brandgray}
        AI Research Engineer \quad $\cdot$ \quad Systems Engineer
        }

    \end{minipage}

\end{center}

\newpage

% ==========================================================
% SLIDE 2: THE CONTEXT
% ==========================================================
\vspace*{0.4in}

\begin{center}
{\fontsize{26}{30}\selectfont\color{brandnavy}\bfseries 
The Initial Setup}
\end{center}

\vspace{0.3in}

{\large\color{brandgray}
Everything looked flawless in the development environment.
}

\vspace{0.15in}

\begin{darkbox}
{\large\textbf{What I Built}}

\vspace{0.05in}
\begin{itemize}[leftmargin=*, itemsep=0.08in, label={\color{brandblue}\textbullet}]
    \item A classification model trained on a cleaned CSV.
    \item Standard in-memory preprocessing pipeline.
    \item A small, efficient data augmentation flow.
    \item Trained locally on a single dev GPU.
\end{itemize}
\end{darkbox}

\vspace{0.2in}

On my laptop, the training loop converged beautifully. Accuracy was stable, validation loss looked great, and inference latency was well within acceptable limits.

\vspace{0.15in}

\begin{infobox}
\textbf{The Fatal Assumption:}\\
I assumed the exact same code and configuration would behave the identically when scaled into a production environment.
\end{infobox}

\vfill
\newpage

% ==========================================================
% SLIDE 3: THE MISTAKE
% ==========================================================
\section*{The Mistake}

\vspace{0.15in}

My biggest error was relying entirely on Jupyter notebook shortcuts and local environmental assumptions.

\vspace{0.2in}

\begin{minipage}[t]{0.45\textwidth}
    {\large\color{brandred}\textbf{What I Did}}

    \begin{itemize}[leftmargin=*, itemsep=0.1in, label={\color{brandred}\textbullet}]
        \item Used implicit library defaults in standard Pandas/NumPy setups.
        \item Tested end-to-end flows manually in cells.
        \item Kept training and serving logic separate.
    \end{itemize}

\end{minipage}
\hfill
\begin{minipage}[t]{0.45\textwidth}

    {\large\color{brandgreen}\textbf{What I Should Have Done}}

    \begin{itemize}[leftmargin=*, itemsep=0.1in, label={\color{brandgreen}\textbullet}]
        \item Enforced strict parity between training and serving.
        \item Pinned all random seeds and default parameters explicitly.
        \item Unified the preprocessing logic.
    \end{itemize}

\end{minipage}

\vspace{0.3in}

\begin{darkbox}

{\large\textbf{The Reality Check}}

\vspace{0.05in}

When it hit production, it silently failed. 

The app returned garbage predictions for a subset of users, costing us hours of frantic debugging and creating a massive spike in our GPU billing.

\end{darkbox}

\vfill
\newpage

% ==========================================================
% SLIDE 4: HOW IT SHOWED UP
% ==========================================================

\section*{How It Manifested}

\vspace{0.15in}

The model didn't crash, it did something much worse. It produced silent errors that looked mathematically valid but were entirely wrong in context.

\vspace{0.15in}

\begin{infobox}[top=10pt, bottom=10pt, left=15pt, right=15pt]
{\large\color{brandblue}\textbf{1. Shifting Feature Distributions}}

\vspace{0.05in}
In production, feature normalization ended up utilizing a completely different code path. The default handler for missing values differed, instantly shifting the distributions of the incoming data.

\vspace{0.05in}
{\small\color{brandgray}\textit{Example: Training used \texttt{fillna(mean)} $\rightarrow$ [2,4,6] normalized (subtract mean, divide by std) to [-1,0,1]. Serving used \texttt{fillna(0)} $\rightarrow$ [2,0,6] normalized to [-0.27,-1.07,1.34]. Same formula, different input geometry.}}
\end{infobox}

\vspace{0.1in}

\begin{infobox}[top=10pt, bottom=10pt, left=15pt, right=15pt]
{\large\color{brandblue}\textbf{2. Silent Precision Differences}}

\vspace{0.05in}
Production batch sizes and data sharding drastically changed runtime memory behavior. Under mixed precision, this caused silent numerical precision differences that altered outputs.

\vspace{0.05in}
{\small\color{brandgray}\textit{Example: FP32 computes 0.1 + 0.2 = 0.30000000000000004 (above threshold 0.3). FP16 stores 0.1 $\approx$ 0.0996 and 0.2 $\approx$ 0.2002, sum = 0.2998 (below threshold). The decision flips silently.}}
\end{infobox}

\vspace{0.1in}

\begin{infobox}[top=10pt, bottom=10pt, left=15pt, right=15pt]
{\large\color{brandblue}\textbf{3. Inconsistent Warmups}}

\vspace{0.05in}
Checkpoint restores in the serving environment loaded a different optimizer state. The model produced highly inconsistent predictions immediately following the warmup phase.

\vspace{0.05in}
{\small\color{brandgray}\textit{Example: Serving restored only model weights, not optimizer buffers. First inference calls used zero momentum $\rightarrow$ predictions deviated until internal state stabilized.}}
\end{infobox}

\vfill
\newpage

% ==========================================================
% SLIDE 5: ROOT CAUSES
% ==========================================================
\section*{Root Cause Analysis}

\vspace{0.05in}

{\fontsize{11}{13}\selectfont\color{brandgray}
A failure of this magnitude is rarely a single bug. It is a compounding of technical and process flaws.
}

\vspace{0.2in}

\begin{minipage}[t]{0.47\textwidth}

{\large\color{brandorange}\textbf{Technical Causes}}

\vspace{0.1in}
\begin{itemize}[leftmargin=*, itemsep=0.15in, label={\color{brandblue}\textbullet}]
    \item \textbf{Path Mismatch:} Preprocessing pipelines differed between the training script and the serving API.
    \item \textbf{Hidden Defaults:} Underlying libraries behaved differently across OS and container environments.
    \item \textbf{Untested Scale:} Mixed precision and checkpointing were never tested end-to-end against production-sized batches.
\end{itemize}

\end{minipage}
\hfill
\begin{minipage}[t]{0.47\textwidth}

{\large\color{brandorange}\textbf{Process Causes}}

\vspace{0.1in}
\begin{itemize}[leftmargin=*, itemsep=0.15in, label={\color{brandblue}\textbullet}]
    \item \textbf{No Staging:} Lacked an environment that mirrored the exact data throughput of production.
    \item \textbf{No Assertions:} Missing unit tests verifying preprocessing parity between raw data and tensors.
    \item \textbf{Notebook Reliance:} Overreliance on interactive experiments without locking in reproducible artifacts.
\end{itemize}

\end{minipage}

\vfill
\newpage

% ==========================================================
% SLIDE 6: WHAT I CHANGED
% ==========================================================

\section*{What I Changed}

\vspace{0.1in}

I implemented three practical changes that stopped these silent failures and completely prevented future regressions.

\vspace{0.15in}

\begin{darkbox}
\textbf{1. Enforce Preprocessing Parity}

Move preprocessing into a single, versioned library used by \textit{both} training and serving. Add unit tests that assert identical outputs for the same raw input dictionaries.
\end{darkbox}

\vspace{0.15in}

\begin{darkbox}
\textbf{2. Staging that Mirrors Production}

Run automated staging jobs with production-like batch sizes and mixed precision explicitly enabled. Validate checkpoint restores and inference warmups under identical memory constraints.
\end{darkbox}

\vspace{0.15in}

\begin{darkbox}
\textbf{3. CI and Instrumentation}

Add Continuous Integration checks that run a tiny end-to-end cycle: \texttt{train $\rightarrow$ checkpoint $\rightarrow$ serve}. Add active monitors for feature distribution drift in production.
\end{darkbox}


\vfill
\newpage

% ==========================================================
% SLIDE 7: THE IMPLEMENTATION
% ==========================================================

\section*{The Implementation}

\vspace{0.1in}

These changes turned a one-off panic fix into a repeatable, scalable process that catches parity issues before they ever reach a user.

\vspace{0.15in}

\textbf{1. Asserting Preprocessing Parity (Python)}

\begin{codebox}{Python}
# test_preprocessing_parity.py
import json
from my_preproc import preprocess_train, preprocess_serve

sample_raw = {
    "text": "Example input",
    "num_feature": None,
    "cat_feature": "A"
}

train_out = preprocess_train(sample_raw)
serve_out = preprocess_serve(sample_raw)

assert train_out == serve_out, f"Mismatch {train_out} != {serve_out}"
print("Preprocessing parity OK")
\end{codebox}

\vspace{0.15in}

\textbf{2. Mirroring Production in Staging (Bash)}

\begin{codebox}{bash}
# run_staging.sh
python train.py --config configs/production_like.yaml \
    --batch-size 256 --mixed-precision true --max-steps 100

python serve.py --checkpoint checkpoints/latest.pt \
    --batch-size 256 --warmup-steps 50
\end{codebox}

\vfill
\newpage

% ==========================================================
% SLIDE 7: THE BIG TAKEAWAY
% ==========================================================

\vspace*{0.35in}

\begin{center}

{\fontsize{14}{16}\selectfont\color{brandblue}\textbf{THE MLOPS TAKEAWAY}}

\vspace{0.3in}

\begin{minipage}{0.88\textwidth}
\centering

{\color{brandblue}\rule{6cm}{1pt}}

\vspace{0.2in}

{\fontsize{24}{28}\selectfont\bfseries\color{brandnavy}
Production parity beats\\
perfect notebook metrics.
}

\vspace{0.15in}

{\fontsize{12}{14}\selectfont\color{brandgray}
If your training and serving environments are not identical in preprocessing, defaults, and runtime behavior, your notebook wins and your users lose. 

\vspace{0.1in}
Ship reproducible pipelines, run small staging jobs that mirror production, and automate your parity checks.
}

\vspace{0.2in}

{\color{brandblue}\rule{6cm}{1pt}}

\vspace{0.3in}

{\fontsize{14}{16}\selectfont\bfseries\color{brandnavy}
What is your worst notebook-to-production surprise?
}

\end{minipage}

\end{center}

\vfill

% Full footer
\begin{center}

\rule{0.6\textwidth}{0.5pt}

\vspace{0.12in}

{\fontsize{11}{13}\selectfont\color{brandgray}
Building AI systems by understanding the mathematics,\\
engineering, and intuition behind every architecture.
}

\vspace{0.12in}

{\fontsize{10}{12}\selectfont\color{brandgray}
\begin{tabular}{@{}l l}
GitHub: & \texttt{github.com/MrHeaven1y} \\
Portfolio: & \texttt{mrheaven1y.github.io} \\
LinkedIn: & \texttt{linkedin.com/in/dibayendu-mukherjee-bb897b267}
\end{tabular}
}

\vspace{0.12in}

{\fontsize{14}{16}\selectfont\color{brandnavy}\textbf{Dibayendu Mukherjee}}

\end{center}

\vspace*{0.2in}

\newpage

% ==========================================================
% SLIDE 8: COMING UP NEXT
% ==========================================================

\vspace*{0.5in}

\begin{center}
    {\fontsize{14}{16}\selectfont\color{brandblue}\textbf{COMING UP NEXT...}}
    
    \vspace{0.4in}
    
    {\fontsize{32}{38}\selectfont\color{brandnavy}\bfseries
    Why Your Training Loss Looks Great\\[0.15in]
    but Your Model Still Fails in Production
    }
    
    \vspace{0.4in}
    
    \begin{minipage}{0.7\textwidth}
        \centering
        \color{brandgray}\fontsize{14}{18}\selectfont
        Evaluation gaps, data leakage, and the silent drift\\
        between validation and reality.
    \end{minipage}
\end{center}

\vfill

\begin{center}
    \rule{0.6\textwidth}{0.5pt}
    
    \vspace{0.15in}
    
    {\fontsize{11}{13}\selectfont\color{brandgray}
    More production-grade MLOps deep dives coming soon.
    }
    
    \vspace{0.12in}
    
    {\fontsize{14}{16}\selectfont\color{brandnavy}\textbf{Dibayendu Mukherjee}}
\end{center}

\vspace*{0.2in}

\newpage